% ============================================================================
%  MATEMÁTICAS OPERATIVAS — Sesión 13
%  Semana 13: 2 al 7 de noviembre de 2026  (lunes 2: festivo, Todos los Santos)
%  3.1 Expresiones algebraicas y operaciones con polinomios
%  3.2 Productos notables
%  Duración prevista: 3 horas
% ============================================================================
\documentclass[aspectratio=169,10pt]{beamer}

\usepackage[utf8]{inputenc}
\IfFileExists{spanish.ldf}{\usepackage[spanish,es-noshorthands]{babel}}{}
\usepackage{amsmath,amssymb}
\usepackage{booktabs}
\usepackage{tikz}

\usetheme{Uniminuto}

\uniasignatura{Matemáticas Operativas}
\uniprofesor{Andrés Felipe Puerta Velez}
\unicorreo{andres.puerta-v@uniminuto.edu.co}
\uninrc{CBASBL011}
\unisemana{Sesión 13 — Semana 13 · 2 al 7 de noviembre de 2026}

\newcommand{\listatiny}{\setbeamerfont{itemize/enumerate body}{size=\tiny}\tiny}

\begin{document}

\portada

\begin{frame}{Orden del día}
  \begin{enumerate}
    \item Retroalimentación del Examen 2
    \item \azul{Tema 3.1: expresiones algebraicas y polinomios}
    \item Sumar, restar y multiplicar polinomios
    \item \azul{Tema 3.2: los cinco productos notables}
    \item Aplicaciones y cálculo mental
    \item Taller principal, errores frecuentes y cierre
  \end{enumerate}

  \vspace{0.25cm}
  {\footnotesize\color{uniGris} Sesión de 3 horas.
  \destacar{Lunes 2 fue festivo} (Todos los Santos). Empieza la Unidad 3: los
  mismos procedimientos de siempre, pero con letras. \azul{Nada de lo que
  viene es una operación nueva.}}
\end{frame}

% ===========================================================================
\seccion{3.1 Expresiones algebraicas y polinomios}
% ===========================================================================

\begin{frame}{Vocabulario mínimo}
  \scriptsize
  \begin{block}{Las piezas}
    \tiny
    Un \azul{término} es un producto de números y letras: $-5x^{2}y$. Su
    \azul{coeficiente} es $-5$; su \azul{parte literal}, $x^{2}y$; su
    \azul{grado}, la suma de exponentes ($3$). Un \azul{polinomio} es una suma
    de términos.
  \end{block}

  \vspace{0.1cm}
  \begin{columns}[T,onlytextwidth]
    \begin{column}{0.48\textwidth}
      \begin{exampleblock}{Términos semejantes}
        \tiny
        Tienen \destacar{la misma parte literal}. Solo esos se suman:
        \[ 7x^{2} + 3x^{2} = 10x^{2} \]
        \[ 7x^{2} + 3x^{3} \ \text{no se reduce} \]
        Es exactamente la regla de los radicales semejantes de la sesión 8.
      \end{exampleblock}
    \end{column}
    \begin{column}{0.48\textwidth}
      \begin{alertblock}{Por qué funciona}
        \tiny
        Por la propiedad \azul{distributiva} de la sesión 1:
        \[ 7x^{2}+3x^{2} = (7+3)x^{2} \]
        \destacar{Sumar términos semejantes es factorizar la parte literal.}
        Con partes literales distintas no hay nada común que sacar.
      \end{alertblock}
    \end{column}
  \end{columns}

  \vspace{0.08cm}
  {\tiny\color{uniGris} Prueba rápida: con $x=2$, $7x^{2}+3x^{2}=40$ y
  $10x^{2}=40$ \checkmark; en cambio $7x^{2}+3x^{3}=52$, que no es
  $10x^{5}=320$. \azul{El contraejemplo numérico funciona igual que con
  logaritmos.}}
\end{frame}

\begin{frame}{Sumar, restar y multiplicar polinomios}
  \scriptsize
  \begin{columns}[T,onlytextwidth]
    \begin{column}{0.48\textwidth}
      \begin{exampleblock}{Suma y resta}
        \tiny
        Se agrupan los semejantes. \destacar{Al restar, el signo entra a todo
        el paréntesis:}
        \[ (4x^{3}-2x+1)-(x^{3}-5x^{2}+3) \]
        \[ = 4x^{3}-2x+1-x^{3}+5x^{2}-3 \]
        \[ = \destacar{3x^{3}+5x^{2}-2x-2} \]
      \end{exampleblock}
    \end{column}
    \begin{column}{0.48\textwidth}
      \begin{exampleblock}{Multiplicación}
        \tiny
        Cada término por cada término (distributiva):
        \[ (2x-3)(x+4) \]
        \[ = 2x^{2}+8x-3x-12 \]
        \[ = \destacar{2x^{2}+5x-12} \]
      \end{exampleblock}
    \end{column}
  \end{columns}

  \vspace{0.12cm}
  \begin{block}{Con más términos, el mismo método}
    \tiny
    \[ (x^{2}+2x-1)(x-3)
     = x^{3}-3x^{2}+2x^{2}-6x-x+3
     = \destacar{x^{3}-x^{2}-7x+3} \]
    \azul{Control:} el grado del producto es la \destacar{suma} de los grados
    ($2+1=3$) \checkmark \ Y el número de productos parciales es
    $3\times 2 = 6$: si le salieron menos, se saltó alguno.
  \end{block}

  \vspace{0.05cm}
  {\tiny\color{uniGris} \destacar{El error que más cuesta:} olvidar cambiar
  \azul{todos} los signos al restar. $-(x^{3}-5x^{2}+3)$ es
  $-x^{3}+5x^{2}-3$, no $-x^{3}-5x^{2}+3$.}
\end{frame}

% ===========================================================================
\seccion{3.2 Productos notables}
% ===========================================================================

\begin{frame}{Los cinco productos notables}
  \scriptsize
  \begin{center}
  \begin{tabular}{@{}p{3.6cm}p{4.2cm}p{4.4cm}@{}}
    \toprule
    \textbf{Nombre} & \textbf{Fórmula} & \textbf{Ejemplo} \\
    \midrule
    Cuadrado de una suma &
      $(a+b)^{2} = a^{2}+2ab+b^{2}$ &
      $(x+7)^{2} = x^{2}+14x+49$ \\
    \addlinespace
    Cuadrado de una resta &
      $(a-b)^{2} = a^{2}-2ab+b^{2}$ &
      $(3x-4)^{2} = 9x^{2}-24x+16$ \\
    \addlinespace
    Suma por diferencia &
      $(a+b)(a-b) = a^{2}-b^{2}$ &
      $(x+9)(x-9) = x^{2}-81$ \\
    \addlinespace
    Cubo de un binomio &
      $(a+b)^{3} = a^{3}+3a^{2}b+3ab^{2}+b^{3}$ &
      $(x+2)^{3} = x^{3}+6x^{2}+12x+8$ \\
    \addlinespace
    Binomios con término común &
      $(x+m)(x+n) = x^{2}+(m+n)x+mn$ &
      $(x+5)(x+3) = x^{2}+8x+15$ \\
    \bottomrule
  \end{tabular}
  \end{center}

  \vspace{0.06cm}
  {\tiny\color{uniGris} \destacar{No son fórmulas nuevas:} todas salen de
  multiplicar término a término. Se memorizan porque \azul{ahorran tiempo} y,
  sobre todo, porque en la sesión 14 hay que reconocerlas \destacar{al revés}
  para factorizar.}
\end{frame}

\begin{frame}{El error más caro del álgebra}
  \scriptsize
  \begin{columns}[T,onlytextwidth]
    \begin{column}{0.30\textwidth}
      \begin{alertblock}{Mal}
        \tiny
        $(a+b)^{2} \neq a^{2}+b^{2}$

        $(x+3)^{2} \neq x^{2}+9$

        \vspace{0.04cm}
        \destacar{Falta el doble producto $2ab$.}
      \end{alertblock}
      \begin{exampleblock}{La prueba con $x=1$}
        \tiny
        $(1+3)^{2} = 16$ \quad frente a \quad $1^{2}+3^{2} = 10$.
        \azul{Un solo número basta} para descartar la igualdad.
      \end{exampleblock}
    \end{column}
    \begin{column}{0.30\textwidth}
      \vspace{0.1cm}
      \begin{center}
      \begin{tikzpicture}[scale=0.52,every node/.style={font=\tiny}]
        \draw[uniAzul,thick,fill=uniAzul!14] (0,0) rectangle (2.6,2.6);
        \draw[uniTitulo,thick,fill=uniCrema!60] (2.6,0) rectangle (4.0,2.6);
        \draw[uniTitulo,thick,fill=uniCrema!60] (0,2.6) rectangle (2.6,4.0);
        \draw[uniNaranja,thick,fill=uniNaranja!22] (2.6,2.6) rectangle (4.0,4.0);
        \node at (1.3,1.3) {$a^{2}$};
        \node at (3.3,1.3) {$ab$};
        \node at (1.3,3.3) {$ab$};
        \node at (3.3,3.3) {$b^{2}$};
        \node[below] at (1.3,-0.05) {$a$}; \node[below] at (3.3,-0.05) {$b$};
      \end{tikzpicture}
      \end{center}
    \end{column}
    \begin{column}{0.36\textwidth}
      \begin{block}{Por qué aparece el $2ab$}
        \tiny
        Un cuadrado de lado $a+b$ se parte en cuatro piezas: el cuadrado
        $a^{2}$, \destacar{dos} rectángulos iguales de área $ab$, y el
        cuadradito $b^{2}$.

        \vspace{0.08cm}
        \[ (a+b)^{2} = a^{2}+\destacar{2ab}+b^{2} \]

        \vspace{0.04cm}
        \azul{Los dos rectángulos del medio son justamente el término que se
        olvida.} Quien vea este dibujo una vez no vuelve a cometer el error.
      \end{block}
    \end{column}
  \end{columns}
\end{frame}

\begin{frame}{Ejercicio 1 — productos notables}
  \scriptsize
  \begin{alertblock}{Individual, 10 minutos. Sin desarrollar término a término.}
    \tiny
    \textbf{1.} $(2a+5b)^{2}$ \qquad
    \textbf{2.} $(4x-3y)(4x+3y)$ \qquad
    \textbf{3.} $(2x-1)^{3}$ \qquad
    \textbf{4.} $(x-7)(x+2)$ \qquad
    \textbf{5.} $(x^{2}+3)^{2}$
  \end{alertblock}

  \vspace{0.06cm}
  \begin{block}{Solución}
    \tiny
    \textbf{1.} $\destacar{4a^{2}+20ab+25b^{2}}$ \qquad
    \textbf{2.} $\destacar{16x^{2}-9y^{2}}$ \qquad
    \textbf{3.} $\destacar{8x^{3}-12x^{2}+6x-1}$

    \vspace{0.08cm}
    \textbf{4.} $x^{2}+(-7+2)x+(-7)(2)
    = \destacar{x^{2}-5x-14}$ \qquad
    \textbf{5.} $\destacar{x^{4}+6x^{2}+9}$
  \end{block}

  \vspace{0.05cm}
  {\tiny\color{uniGris} En el \destacar{3}, los signos alternan porque el
  binomio es una resta: $+,-,+,-$. En el \destacar{5}, el ``$a$'' es $x^{2}$,
  así que $a^{2}=x^{4}$ y $2ab = 2\cdot x^{2}\cdot 3 = 6x^{2}$.}
\end{frame}

\begin{frame}{Los productos notables sirven para calcular mental}
  \scriptsize
  \begin{columns}[T,onlytextwidth]
    \begin{column}{0.48\textwidth}
      \begin{exampleblock}{Suma por diferencia}
        \tiny
        \[ 103\cdot 97 = (100+3)(100-3) \]
        \[ = 10\,000 - 9 = \destacar{9\,991} \]
      \end{exampleblock}
    \end{column}
    \begin{column}{0.48\textwidth}
      \begin{exampleblock}{Cuadrado de una suma}
        \tiny
        \[ 52^{2} = (50+2)^{2} \]
        \[ = 2\,500 + 200 + 4 = \destacar{2\,704} \]
      \end{exampleblock}
    \end{column}
  \end{columns}

  \vspace{0.12cm}
  \begin{block}{Otros que salen redondos}
    \tiny
    \[ 98^{2} = (100-2)^{2} = 10\,000-400+4 = \destacar{9\,604} \]
    \[ 45\cdot 55 = (50-5)(50+5) = 2\,500-25 = \destacar{2\,475} \]
  \end{block}

  \vspace{0.08cm}
  {\tiny\color{uniGris} \destacar{Cuándo buscar el atajo:} cuando los dos
  números están \azul{a la misma distancia} de una decena o centena redonda
  (suma por diferencia), o cuando uno está \azul{cerca} de una (cuadrado de
  binomio). No siempre aplica, pero cuando aplica ahorra media hoja.}
\end{frame}

\begin{frame}{Ejercicio 2 — taller principal}
  \begin{alertblock}{En grupos de tres, 20 minutos. Se socializa en el tablero.}
    \scriptsize
    \begin{enumerate}\setlength{\itemsep}{0.0cm}
      \item Reduzca: \ $(3x^{2}-5x+2)+(x^{2}+7x-9)$
      \item Multiplique: \ $3x^{2}(2x-5)$
      \item Multiplique: \ $(x^{2}+2x-1)(x-3)$
      \item Desarrolle: \ $(5x+2)^{2}$ \ y \ $(x+4)^{3}$
      \item Calcule mentalmente: \ $206\cdot 194$ \ y \ $71^{2}$
      \item Un terreno cuadrado de lado $x$ se amplía $6$ m por cada lado.
            Escriba el área nueva como producto notable y desarrollada.
    \end{enumerate}
  \end{alertblock}
\end{frame}

\begin{frame}{Ejercicio 2 — solución}
  \scriptsize
  \begin{columns}[T,onlytextwidth]
    \begin{column}{0.49\textwidth}
      \begin{block}{Puntos 1 a 3}
        \tiny
        \textbf{1.} $\destacar{4x^{2}+2x-7}$

        \vspace{0.08cm}
        \textbf{2.} $\destacar{6x^{3}-15x^{2}}$

        \vspace{0.08cm}
        \textbf{3.} $x^{3}-3x^{2}+2x^{2}-6x-x+3
        = \destacar{x^{3}-x^{2}-7x+3}$
      \end{block}
    \end{column}
    \begin{column}{0.49\textwidth}
      \begin{block}{Puntos 4 a 6}
        \tiny
        \textbf{4.} $\destacar{25x^{2}+20x+4}$; \
        $\destacar{x^{3}+12x^{2}+48x+64}$

        \vspace{0.08cm}
        \textbf{5.} $(200+6)(200-6)=40\,000-36
        =\destacar{39\,964}$; \
        $(70+1)^{2}=4\,900+140+1=\destacar{5\,041}$

        \vspace{0.08cm}
        \textbf{6.} $(x+6)^{2}
        = \destacar{x^{2}+12x+36}$
      \end{block}
    \end{column}
  \end{columns}

  \vspace{0.08cm}
  {\tiny\color{uniGris} El punto \destacar{6} le da sentido al $2ab$: el área
  nueva no es ``el cuadrado viejo más $36$''. Son \azul{dos franjas} de
  $6x$ a los lados más la esquinita de $6\times 6$. Es literalmente el dibujo
  de la lámina anterior.}
\end{frame}

\begin{frame}{Errores frecuentes en esta sesión}
  \scriptsize
  \begin{center}
  \begin{tabular}{@{}p{5.0cm}p{7.0cm}@{}}
    \toprule
    \textbf{Error} & \textbf{Lo correcto} \\
    \midrule
    $(a+b)^{2}=a^{2}+b^{2}$ &
      Falta $2ab$. Compruebe con $a=b=1$: $4\neq 2$ \\
    \addlinespace
    $-(x^{3}-5x^{2}+3) = -x^{3}-5x^{2}+3$ &
      El signo entra a \azul{todos}: $-x^{3}+5x^{2}-3$ \\
    \addlinespace
    $7x^{2}+3x^{3}=10x^{5}$ &
      No son semejantes; no se reduce \\
    \addlinespace
    $(3x)^{2}=3x^{2}$ &
      $=9x^{2}$: el exponente afecta al coeficiente \\
    \addlinespace
    $(a-b)^{2}=a^{2}-b^{2}$ &
      Eso es $(a+b)(a-b)$; el cuadrado da $a^{2}-2ab+b^{2}$ \\
    \addlinespace
    Saltarse productos parciales &
      $(3\ \text{términos})\times(2)$ da $6$ productos \\
    \bottomrule
  \end{tabular}
  \end{center}
\end{frame}

\begin{frame}{Cierre de la sesión}
  \scriptsize
  \begin{block}{Lo que hay que llevarse}
    \tiny
    \textbf{1.} Solo se suman términos \azul{semejantes}, y es la distributiva
    aplicada al revés. \quad
    \textbf{2.} Al restar un polinomio, cambian \destacar{todos} los signos.
    \quad
    \textbf{3.} Los cinco productos notables se memorizan porque en la sesión
    14 hay que leerlos \azul{al revés}. \quad
    \textbf{4.} $(a\pm b)^{2}$ siempre lleva el doble producto. \quad
    \textbf{5.} Un contraejemplo numérico descarta cualquier igualdad falsa.
  \end{block}

  \vspace{0.1cm}
  \begin{alertblock}{Para la próxima sesión}
    \tiny
    Resolver los \destacar{20 ejercicios propuestos} y la guía de la unidad en
    Aulas Virtuales.

    \vspace{0.05cm}
    En la sesión 14 vemos \azul{factorización y fracciones algebraicas}: el
    tema más rentable de la unidad. Sin factorizar no se simplifican
    fracciones ni se resuelven ecuaciones. \destacar{Llegue con los cinco
    productos notables sabidos de memoria.}
  \end{alertblock}
\end{frame}

\begin{frame}{Ejercicios propuestos}
  \vspace{-0.1cm}
  \begin{beamercolorbox}[wd=\textwidth,sep=0.10cm]{block body}
    \centering\scriptsize
    \textbf{1--6:} reduzca u opere. \quad \textbf{7--16:} desarrolle. \quad
    \textbf{17--20:} cálculo mental y aplicación.
  \end{beamercolorbox}
  \listatiny
  \vspace{0.15cm}
  \begin{columns}[T,onlytextwidth]
    \begin{column}{0.47\textwidth}
      \begin{enumerate}\setlength{\itemsep}{0.16cm}
        \item $(5x^{2}-3x)+(2x^{2}+8x-4)$
        \item $(7y^{3}+2y)-(3y^{3}-y+6)$
        \item $4x(3x^{2}-2x+5)$
        \item $(x+6)(x-2)$
        \item $(2x-1)(3x+4)$
        \item $(x^{2}-x+2)(x+1)$
        \item $(x+11)^{2}$
        \item $(2x-7)^{2}$
        \item $(3a+4b)^{2}$
        \item $(x+12)(x-12)$
      \end{enumerate}
    \end{column}
    \begin{column}{0.50\textwidth}
      \begin{enumerate}\setlength{\itemsep}{0.16cm}
        \setcounter{enumi}{10}
        \item $(5x-2y)(5x+2y)$
        \item $(x-3)^{3}$
        \item $(2x+1)^{3}$
        \item $(x+8)(x+5)$
        \item $(x-9)(x+4)$
        \item $(x^{3}+2)^{2}$
        \item $104\cdot 96$
        \item $61^{2}$
        \item $99^{2}$
        \item Un cuadrado de lado $y$ se reduce $3$ m por lado: escriba el
              área nueva desarrollada
      \end{enumerate}
    \end{column}
  \end{columns}
\end{frame}

\begin{frame}{Respuestas}{Autoevaluación de los ejercicios propuestos}
  \listatiny
  \vspace{0.15cm}
  \begin{columns}[T,onlytextwidth]
    \begin{column}{0.47\textwidth}
      \begin{enumerate}\setlength{\itemsep}{0.16cm}
        \item $7x^{2}+5x-4$
        \item $4y^{3}+3y-6$
        \item $12x^{3}-8x^{2}+20x$
        \item $x^{2}+4x-12$
        \item $6x^{2}+5x-4$
        \item $x^{3}+x+2$
        \item $x^{2}+22x+121$
        \item $4x^{2}-28x+49$
        \item $9a^{2}+24ab+16b^{2}$
        \item $x^{2}-144$
      \end{enumerate}
    \end{column}
    \begin{column}{0.50\textwidth}
      \begin{enumerate}\setlength{\itemsep}{0.16cm}
        \setcounter{enumi}{10}
        \item $25x^{2}-4y^{2}$
        \item $x^{3}-9x^{2}+27x-27$
        \item $8x^{3}+12x^{2}+6x+1$
        \item $x^{2}+13x+40$
        \item $x^{2}-5x-36$
        \item $x^{6}+4x^{3}+4$
        \item $9\,984$
        \item $3\,721$
        \item $9\,801$
        \item $(y-3)^{2}=y^{2}-6y+9$
      \end{enumerate}
    \end{column}
  \end{columns}
\end{frame}

\begin{frame}{Referencias}
  \footnotesize
  \begin{itemize}
    \item[{[1]}] Ramos del Olmo, Á. M. (2016). \textit{Matemáticas básicas
      para el acceso a la universidad}. Ediciones Pirámide.
      \texttt{https://elibro.net/es/ereader/uniminuto/49134}
    \item[{[2]}] González Pulido, J. W., Chacón Benavides, J. A. \&
      Fonseca Correa, L. Á. (2023). \textit{Matemática básica}
      (1.\textsuperscript{a} ed.). Universidad Mariana.
      \texttt{https://elibro.net/es/ereader/uniminuto/280647}
    \item[{[3]}] Torrijos, M. (2018). \textit{Apuntes de matemática básica}.
      Universidad Católica de Colombia.
      \texttt{https://elibro.net/es/ereader/uniminuto/197048}
    \item[{[4]}] Ramírez, A. \& Rojas, L. (2016). \textit{Matemáticas básicas
      con aplicaciones a la ingeniería}. Ecoe Ediciones.
      \texttt{https://elibro.net/es/ereader/uniminuto/114355}
  \end{itemize}

  \vspace{0.3cm}
  \begin{center}
    {\Large\color{uniAzul}\textbf{¿Preguntas?}}
  \end{center}
\end{frame}

\end{document}
